\documentclass{article}
\usepackage[utf8]{inputenc}
\usepackage[vietnamese]{babel} 
\usepackage{times} % Font chữ Times New Roman

% --- CẤU HÌNH KHOẢNG CÁCH DÒNG 1.5 ĐỀU NHAU ---
\usepackage{setspace}
\onehalfspacing % Ép toàn bộ văn bản giãn dòng 1.5 đều nhau

% --- CẤU HÌNH THỤT ĐẦU DÒNG CHUẨN ---
\usepackage{indentfirst} % Ép đoạn văn đầu tiên sau tiêu đề phải thụt đầu dòng
\setlength{\parindent}{0.5in} % Độ rộng thụt vào đầu dòng (khoảng 1.27cm)

\usepackage{titlesec}
\titlelabel{\thetitle.\quad} % Tự động thêm dấu chấm

% Tối ưu lề
\setlength{\oddsidemargin}{0in}
\setlength{\evensidemargin}{0in}
\setlength{\textwidth}{6.5in}
\setlength{\topmargin}{-0.5in}
\setlength{\textheight}{9in}

% Thư viện vẽ hình TikZ để dựng sơ đồ
\usepackage{tikz}
\usetikzlibrary{shapes.geometric, arrows, positioning}

\begin{document}
	
	\section{MÔ HÌNH HÓA DỮ LIỆU MỨC KHÁI NIỆM (CONCEPTUAL ERD)}
	
	\subsection{Khái Niệm, Quy Tắc Nghiệp Vụ và Sơ Đồ Chi Tiết}
	
	% --- PHẦN 1: ĐỊNH NGHĨA THỰC THỂ ---
	\noindent \textbf{1. Xác định chi tiết các thực thể tham gia hệ thống.} \\
	\indent Dựa trên các yêu cầu đặc tả chức năng từ hệ thống định hướng nghề nghiệp cá nhân hóa, cấu trúc dữ liệu ở mức khái niệm bao gồm 6 thực thể cốt lõi sau:
	\begin{itemize}
		\item \textbf{Student Profile (Hồ sơ Sinh viên):} Thực thể trung tâm lưu trữ dữ liệu định danh của sinh viên (MSSV, Họ tên, Lớp). Đồng thời, thực thể này chịu trách nhiệm tích hợp kết quả học tập (Transcript) và Token kết nối tài khoản GitHub phục vụ cho AI phân tích kỹ năng.
		\item \textbf{Tech Path (Lộ trình Công nghệ):} Định nghĩa các vị trí nghề nghiệp tiêu chuẩn trong ngành Kỹ thuật Phần mềm (ví dụ: Frontend Developer, Backend Developer, DevOps Engineer, Data Engineer). Mỗi lộ trình đóng vai trò là cái khung định hướng tổng quát.
		\item \textbf{Skill Node (Nút Kỹ năng):} Các đơn vị kiến thức hoặc công nghệ cấu thành nên một Tech Path (ví dụ: Java, Docker, Git, SQL). Các nút này được liên kết phân cấp và định trọng số ưu tiên theo thứ tự từ cơ bản đến nâng cao.
		\item \textbf{Course Repo (Kho Tài nguyên):} Hệ thống lưu trữ các liên kết bài học, tài liệu kỹ thuật chuẩn, hoặc các khóa học mở (MOOCs) tương ứng trực tiếp với từng nút kỹ năng để định hướng sinh viên học tập.
		\item \textbf{Job Trend (Xu hướng Việc làm):} Kết quả thu được từ tiến trình tự động cào dữ liệu (Web Scraper) từ các sàn tuyển dụng. Thực thể này lưu trữ các chỉ số về tần suất xuất hiện của từ khóa công nghệ, số lượng tin tuyển dụng đang mở.
		\item \textbf{Mentor Session (Phiên Cố vấn):} Ghi nhận thông tin tương tác lịch sử giữa sinh viên với Cố vấn học tập hoặc Chuyên gia công nghiệp (Thời gian, nội dung câu hỏi, đánh giá từ Mentor, hoặc log chat với AI Mentor).
	\end{itemize}
	
	\vspace{0.3cm}
	
	% --- PHẦN 2: QUY TẮC NGHIỆP VỤ ---
	\noindent \textbf{2. Quy tắc Nghiệp vụ Giữa Các Thực Thể (Business Logic).} \\
	\indent Mối quan hệ và ràng buộc dữ liệu mức khái niệm được thiết lập dựa trên các quy tắc nghiệp vụ sau:
	\begin{enumerate}
		\item Một sinh viên (\textbf{Student Profile}) có thể chủ động đăng ký học nhiều lộ trình nghề nghiệp khác nhau, và một lộ trình (\textbf{Tech Path}) cũng được theo đuổi bởi nhiều sinh viên khác nhau $\rightarrow$ \textbf{Quan hệ Nhiều - Nhiều ($N-M$)}.
		\item Một lộ trình (\textbf{Tech Path}) được cấu thành từ tập hợp nhiều kỹ năng, và một kỹ năng (\textbf{Skill Node}) có thể dùng chung cho nhiều lộ trình nghề nghiệp $\rightarrow$ \textbf{Quan hệ Nhiều - Nhiều ($N-M$)}.
		\item Để đảm bảo tính đa dạng của tài liệu, một kỹ năng (\textbf{Skill Node}) phải được cung cấp và liên kết với nhiều tài nguyên học tập (\textbf{Course Repo}) $\rightarrow$ \textbf{Quan hệ Một - Nhiều ($1-N$)}.
		\item Trong quá trình định hướng, một sinh viên (\textbf{Student Profile}) có thể đặt lịch và tham gia vào nhiều phiên hỗ trợ (\textbf{Mentor Session}) khác nhau $\rightarrow$ \textbf{Quan hệ Một - Nhiều ($1-N$)}.
		\item Dữ liệu thị trường (\textbf{Job Trend}) sau khi phân tích từ khóa sẽ dùng để cập nhật trạng thái nhu cầu thực tế cho nhiều kỹ năng (\textbf{Skill Node}) tương ứng $\rightarrow$ \textbf{Quan hệ Nhiều - Nhiều ($N-M$)}.
	\end{enumerate}
	
	\vspace{0.3cm}
	
	% --- PHẦN 3: SƠ ĐỒ ERD ---
	\noindent \textbf{3. Sơ đồ Thực thể Mối quan hệ mức Khái niệm (Conceptual ERD Diagram).} \\
	\indent Sơ đồ dưới đây đã được cấu trúc lại theo dạng mạng lưới (Grid) 3 tầng để tránh hiện tượng chồng lấn các thực thể và đảm bảo tính thẩm mỹ khi xuất bản tài liệu, đồng thời sửa đổi chuẩn xác bản số kết nối giữa Skill Node và Course Repo:
	
	\begin{figure}[htbp]
		\begin{center}
			\begin{singlespace} 
				\begin{tikzpicture}[
					entity/.style={rectangle, draw=blue!70, fill=blue!5, thick, minimum width=3.2cm, minimum height=0.9cm, rounded corners=2pt},
					relationship/.style={diamond, draw=red!70, fill=red!5, thick, minimum width=2.2cm, minimum height=0.9cm, aspect=1.5},
					every node/.style={font=\small\bfseries},
					node distance=1.4cm and 1.2cm
					]
					
					% --- TẦNG 1: QUẢN LÝ ĐỊNH HƯỚNG ---
					\node[entity] (student) {Student Profile};
					\node[relationship] (select) [right=of student] {Chọn};
					\node[entity] (path) [right=of select] {Tech Path};
					
					% --- TẦNG 2: ĐƯỜNG KẾT NỐI TRUNG GIAN ---
					\node[relationship] (have) [below=of student] {Tham gia};
					\node[relationship] (contain) [below=of path] {Chứa};
					
					% --- TẦNG 3: CHI TIẾT ---
					\node[entity] (mentor) [below=of have] {Mentor Session};
					\node[entity] (skill) [below=of contain] {Skill Node};
					
					% --- TẦNG 4: NHÁNH PHÂN TÁN TỪ SKILL NODE  ---
					\node[relationship] (link) [below left=of skill] {Liên kết};
					\node[entity] (course) [below=of link] {Course Repo};
					
					\node[relationship] (update) [below right=of skill] {Cập nhật};
					\node[entity] (job) [below=of update] {Job Trend};
					
					% --- ĐƯỜNG NỐI VÀ GÁN BẢN SỐ (CARDINALITY) ---
					\path[draw, thick] (student) -- node[above, near start] {N} (select);
					\path[draw, thick] (select) -- node[above, near end] {M} (path);
					
					\path[draw, thick] (student) -- node[left, near start] {1} (have);
					\path[draw, thick] (have) -- node[left, near end] {N} (mentor);
					
					\path[draw, thick] (path) -- node[left, near start] {N} (contain);
					\path[draw, thick] (contain) -- node[left, near end] {M} (skill);
					
					%1 Skill có N Course
					\path[draw, thick] (skill) -- node[above left, near start] {1} (link);
					\path[draw, thick] (link) -- node[below left, near end] {N} (course);
					
					\path[draw, thick] (skill) -- node[above right, near start] {N} (update);
					\path[draw, thick] (update) -- node[below right, near end] {M} (job);
					
				\end{tikzpicture}
			\end{singlespace}
		\end{center}
	
		\caption{Sơ đồ ERD mức Khái niệm hệ thống Định hướng \& Lộ trình học tập}
		\label{fig:fixed_conceptual_erd}
	\end{figure}
	
\end{document}